% =============================================================================
\section{Method of Solution}
\label{sec:method}
% =============================================================================

% -----------------------------------------------------------------------------
\subsection{System Structure}
\label{sec:method:structure}
% -----------------------------------------------------------------------------

At all operating configurations, \cref{eq:mass-balance,eq:energy-balance,eq:solid-temp}
form a pure system of ODEs:
\begin{equation}
  \dot{\vect{x}} = \vect{f}(\vect{x}, t).
\end{equation}
When one or more tanks enter Configuration~B, an algebraically computed
compensating heat flux $\dot{Q}_{HX,i}$ (\cref{sec:dae:configB}) is added to
the temperature ODE right-hand side at each time step.  Because
$\dot{Q}_{HX,i}$ is an explicit function of the current state $(T_i, \rho_i)$
and the prescribed mission discharge rate, no implicit equations arise and the
system remains a pure ODE throughout.

% -----------------------------------------------------------------------------
\subsection{Per-Step Heat-flux Evaluation for Configuration~B}
\label{sec:method:configB-solve}
% -----------------------------------------------------------------------------

\Cref{alg:rhs} summarises the right-hand-side evaluation procedure.

\begin{figure}[ht]
\begin{tikzpicture}[
    box/.style    = {draw, rounded corners=3pt, fill=blue!5, minimum width=8cm,
                     minimum height=0.7cm, align=left, font=\small, inner sep=6pt},
    diamond/.style= {draw, diamond, fill=yellow!15, minimum width=2.5cm,
                     minimum height=0.7cm, align=center, font=\small, inner sep=4pt},
    arr/.style    = {-{Latex[length=5pt]}, thick},
    node distance = 0.35cm,
  ]
  \node[box] (s1)  {\textbf{1.} Reconstruct $\vect{x}_i = (m_i, T_i, T_{s,i})$ from state vector $\vect{y}$.};
  \node[box] (s2) [below=of s1]
                   {\textbf{2.} Compute all coupling edge flows $\mdot_e^{\text{coup}}$.};
  \node[box] (s3) [below=of s2]
                   {\textbf{3.} For each tank~$i$: determine configuration (A/B/C).};
  \node[box] (s4) [below=of s3]
                   {\textbf{4.\ Config~B tanks:} evaluate compensating heat flux
                    $\dot{Q}_{HX,i}$ via \cref{eq:config-b-qhx};
                    add to energy balance RHS.};
  \node[box] (s5) [below=of s4]
                   {\textbf{5.\ Config~C tanks:} compute vent edge flow $\mdot_e^{\text{vent}}$
                    via \cref{eq:vent-rate}.};
  \node[box] (s6) [below=of s5]
                   {\textbf{6.} Compute thermal model: $\dot{Q}_{s,i}$, $\dot{T}_{s,i}$.};
  \node[box] (s7) [below=of s6]
                   {\textbf{7.} Assemble $\dot{m}_i$, $\dot{T}_i$, $\dot{T}_{s,i}$
                    from \cref{eq:mass-balance}, \cref{eq:energy-balance}, and \cref{eq:solid-temp}.};
  \node[box] (s8) [below=of s7]
                   {\textbf{8.} Return $\dot{\vect{y}}$.};

  \foreach \a/\b in {s1/s2, s2/s3, s3/s4, s4/s5, s5/s6, s6/s7, s7/s8}
    \draw[arr] (\a.south) -- (\b.north);
\end{tikzpicture}
\caption{RHS evaluation procedure for one ODE time step.  Step~4 algebraically
evaluates $\dot{Q}_{HX,i}$ for any Config-B tanks and injects it into the
energy balance; step~5 handles venting.  All remaining steps are standard ODE
evaluations.}
\label{alg:rhs}
\end{figure}

Steps 2--5 constitute the explicit algebraic preprocessing.  The coupling flows
(step~2) use the current state and are treated as prescribed inputs to the
remainder of the RHS; no fixed-point iteration is required because the valve
models and the Config-B heat flux are all explicit functions of the state.

% -----------------------------------------------------------------------------
\subsection{Time Integration}
\label{sec:method:integration}
% -----------------------------------------------------------------------------

The ODE system (with Config-B heat fluxes included in the RHS where applicable)
is integrated using the \textsc{lsoda} solver from \texttt{scipy.integrate},
accessed through the \texttt{solve\_ivp} interface.

\textsc{lsoda} was chosen for the following reasons:

\begin{itemize}[nosep, left=0pt]
  \item It switches automatically between Adams (non-stiff) and BDF (stiff)
        formulæ, which is important because stiffness changes throughout the
        mission as configurations switch and coupling valves open or close.
  \item Relative and absolute tolerances are set by the user (\texttt{rtol},
        \texttt{atol}) and are tightened by default to \num{1e-4} and
        \num{1e-7} respectively to control mass-conservation drift.
  \item It produces dense output, enabling post-processing at arbitrary time
        points without re-running the simulation.
\end{itemize}

The integrator step size is fully adaptive; no explicit step-size control is
needed from user code.

% -----------------------------------------------------------------------------
\subsection{Stopping Criteria}
\label{sec:method:stopping}
% -----------------------------------------------------------------------------

Integration terminates when any of the following events is detected:

\begin{center}
\begin{tabularx}{\linewidth}{@{} l l X @{}}
\toprule
\textbf{Criterion} & \textbf{Variable} & \textbf{Description} \\
\midrule
Minimum density  & $\rho_i \le \rho_{\min}$ &
  The tank is effectively empty for the given mission.
  Default: $\rho_{\min} = \SI{5.8}{\kilogram\per\cubic\metre}$ (gaseous H$_2$). \\[4pt]
Target density   & $\rho_i \ge \rho_{\text{target}}$ &
  Used in refuel missions.  Integration stops when the tank has been filled
  to the desired density. \\[4pt]
Duration         & $t \ge t_{\text{mission}}$ &
  Hard time limit derived from the mission profile duration. \\
\bottomrule
\end{tabularx}
\end{center}

All criteria are checked at each RHS evaluation.  When a criterion is met, a
flag is set and the RHS returns a near-zero vector, allowing the solver to
cleanly reach the termination point without an abrupt step.

% -----------------------------------------------------------------------------
\subsection{Configuration Switching and Hysteresis}
\label{sec:method:hysteresis}
% -----------------------------------------------------------------------------

Configuration detection must be numerically robust to avoid chattering near
transition boundaries.  The following hysteresis scheme is used for Config~B:

\begin{align}
  \text{A} \to \text{B} &: \quad P_i \le P_{\min,i}(1 - \delta), \\
  \text{B} \to \text{A} &: \quad P_i >  P_{\min,i}(1 + \delta),
\end{align}
with $\delta = 0.01$ (1\%).  The current configuration is stored as instance
state across RHS calls so that the hysteresis condition can reference the
previous configuration.

Config~C (venting) has no hysteresis: it activates immediately when
$P_i \ge P_{\text{vent},i}$ and deactivates when the vent flow brings the
pressure back below the threshold.
